\subsection{可分代数}

定义1. --- 设 A 是域 K 上的一个交换代数。若对 K 的每个扩张 L，环 L ⊗K A 都是既约的，则称 A 在 K 上是可分的，或称为可分 K-代数。

每个可分代数显然是既约的。部分逆命题见定理3（V，第125页）。

例. --- 1) 设 A 是一个多项式代数 \(K[X_{i}]_{i \in I}\) 。对 K 的每个扩张 L，环 \(L \otimes_{K} A\) 同构于 \(L[X_{i}]_{i \in I}\) （III，第449页，注记2），因此是一个整环（IV，第9页，命题8）。换言之，域 K 上的每个多项式代数都是可分 K-代数。

\begin{enumerate}
\def\labelenumi{\arabic{enumi})}
\setcounter{enumi}{1}
\item
  设 A 是域 K 上的一个有限次交换代数。为使 A 是可分的，充要条件是它是平展的（V，第34页，定理4）。
\item
  设 E 是域 K 的一个代数扩张。若 L 是 K 的一个扩张，则环 \(L \otimes_{K} E\) 是子环 \(L \otimes_{K} F\) 的并，其中 F 遍历 E 的所有有限次子扩张；因此环 \(L \otimes_{K} E\) 是既约的当且仅当对 E 的每个在 K 上有限的子扩张 F， \(L \otimes_{K} F\) 都是既约的。顾及例2，我们看到 E 是上述定义1意义下的可分代数，当且仅当它是 V 第36页定义1意义下的可分代数扩张。
\end{enumerate}

命题3. --- 设K是一个域。

\begin{enumerate}
\def\labelenumi{\alph{enumi})}
\item
  可分 K-代数的每个子代数都是可分的。
\item
  可分 K-代数的每个直接极限都是可分的。
\item
  可分 K-代数的每个乘积都是可分的。
\item
  设 A 是一个 K-代数， \(K'\) 是 K 的一个扩张。为使 A 是可分的，充要条件是由 A 经标量扩张得到的 \(K'\) -代数 \(\mathbf{A}_{(\mathbf{K}')}\) 是可分的。
\end{enumerate}

设 \(\mathbf{L}\) 是 \(\mathbf{K}\) 的一个扩张。设 \(\mathbf{A}\) 是 \(\mathbf{K}\) 上的一个可分代数， \(\mathbf{B}\) 是 \(\mathbf{A}\) 的一个子代数；则环 \(\mathbf{L} \otimes_{\mathbf{K}} \mathbf{A}\) 是既约的，且 \(\mathbf{L} \otimes_{\mathbf{K}} \mathbf{B}\) 同构于 \(\mathbf{L} \otimes_{\mathbf{K}} \mathbf{A}\) 的一个子环，因而是既约的。于是 \(\mathbf{B}\) 是可分的，这就证明了 \(a)\) 。同理可证 \(b)\) ，利用 \(L \otimes_{K} \varinjlim A_i\) 与 \(\varinjlim L \otimes_K A_i\) 的典范同构（II，第290页，命题7）；而 \(c)\) 由指出 \(L \otimes_K \left( \prod_{i \in I} A_i \right)\) 同构于 \(\prod_{i \in I} (L \otimes_K A_i)\) 的一个子环（II，第306页，命题15）得到证明。

我们将使用 \(d\) ) 的记号。对每个扩张 \(\mathbf{L}'\) （ \(\mathbf{K}'\) 的），环 \(\mathbf{L}' \otimes_{\mathbf{K}'} \mathbf{A}_{(\mathbf{K}')}\) 与 \(\mathbf{L}' \otimes_{\mathbf{K}} \mathbf{A}\) 是同构的（II，第278页，命题2）。我们得出结论：若 \(\mathbf{A}\) 是可分 \(\mathbf{K}\) -代数，则 \(\mathbf{A}_{(\mathbf{K}')}\) 是可分 \(\mathbf{K}'\) -代数。反之，设 \(\mathbf{A}_{(\mathbf{K}')}\) 是可分 \(\mathbf{K}'\) -代数；前面的注记表明 \(\mathbf{L}' \otimes_{\mathbf{K}} \mathbf{A}\) 是既约的（对每个扩张 \(\mathbf{L}'\) （ \(\mathbf{K}\) 的、包含 \(\mathbf{K}'\) 作为子扩张））。设 \(\mathbf{L}\) 是 \(\mathbf{K}\) 的一个扩张；由补注（V，第13页），存在一个扩张 \(\mathbf{L}'\) （ \(\mathbf{K}\) 的、包含 \(\mathbf{K}'\) 作为子扩张的）以及一个 \(\mathbf{K}\) -同态，从 \(\mathbf{L}\) 到 \(\mathbf{L}'\) ；于是环 \(\mathbf{L} \otimes_{\mathbf{K}} \mathbf{A}\) 同构于 \(\mathbf{L}' \otimes_{\mathbf{K}} \mathbf{A}\) 的一个子环，因而是既约的。这就证明了 \(\mathbf{A}\) 是可分 \(\mathbf{K}\) -代数。

命题4. --- 设 A 是域 K 上的一个可分代数，并设 B 是 A 的全分式环；则 K-代数 B 是可分的。

设 S 是 A 的可消去元集合。我们已把 A 与 B 的一个子环等同（I，第113页）；此外，S 的每个元素在 B 中可逆，且 B 的每个元素具有形式 \(as^{-1}\) ，其中 \(a \in \mathbf{A}\) ， \(s \in \mathbf{S}\) 。设 \(L\) 是 K 的一个扩张， \(x\) 是 \(\mathbf{L} \otimes_{\mathbf{K}} \mathbf{B}\) 的一个幂零元。则 \(x\) 可写成形式 \(x = \sum_{i=1}^{n} y_i \otimes a_i s_i^{-1}\) ，其中 \(y_i \in \mathbf{L}\) ， \(a_i \in \mathbf{A}\) ， \(s_i \in \mathbf{S}\) （ \(1 \leqslant i \leqslant n\) ）。若令 \(s = s_1 \ldots s_n\)

则 \(x(1 \otimes s)\) 属于子环 \(L \otimes_{K} A\) （ \(L \otimes_{K} B\) 的）；由于 \(x(1 \otimes s)\) 是幂零的且 A 在 K 上可分，我们有 \(x(1 \otimes s) = 0\) ，又由于 s 在 B 中可逆，我们得到 x = 0。这就证明了环 \(L \otimes_{K} B\) 是既约的，命题得证。

命题5. --- 设 K 是一个域，A、B 是交换 K-代数。若 A 是既约的且 B 是可分的，则 A ⊗K B 是既约的。

由命题2的推论（V，第119页），A 同构于一个乘积 \(\prod_{i\in I}L_i\) 的子代数，其中 \(L_{i}\) 对每个 \(i\in I\) 是 K 的一个扩张。因此 \(\mathbf{A}\otimes_{\mathbf{K}}\mathbf{B}\) 同构于

\(\left(\prod_{i\in I}L_i\right)\otimes_K B\) 的一个子环，且后一个环同构于

\(\prod_{i\in I}(\mathbf{L}_i\otimes_{\mathbb{K}}\mathbf{B})\) 的一个子环（II，第306页，命题15）。由于 \(\mathbf{B}\) 是可分的，诸环

\(L_{i} \otimes_{K} B\) 中的每一个都是既约的，因而 \(\prod_{i \in I} (L_{i} \otimes_{K} B)\) 亦然，从而 \(A \otimes_{K} B\) 更是如此。

推论1. --- 设 K 是一个域，L 是 K 的一个可分扩张，f 是 K{[}X{]} 中的一个多项式。若 f 在 K{[}X{]} 中没有重因子，则它在 L{[}X{]} 中也没有重因子。

若 \(f\) 在 \(\mathbf{K}[\mathbf{X}]\) 中没有重因子，则商环 \(\mathbf{K}[\mathbf{X}] / (f)\) 是既约的；因为 \(f\) 是两两互素的不可约多项式 \(f_i\) 的乘积，因此

K{[}X{]}/(f) 由 I 第110页命题9同构于域的乘积 K{[}X{]}/(f\_i)。由命题5，环 L{[}X{]}/(f) （它同构于 L ⊗\_K K{[}X{]}/(f) ）是既约的；若 g 是 L{[}X{]} 的一个使得 g\^{}2 整除 f 的非常数多项式，则 fg\^{}\{-1\} 在 L{[}X{]}/(f) 中的剩余类是一个非零幂零元；因此 f 在 L{[}X{]} 中没有重因子。

推论2. --- 设 A 和 B 是两个交换 K-代数。若 A 和 B 都是可分的，则 \(\mathbf{A} \otimes_{\mathbf{K}} \mathbf{B}\) 亦然。

设 L 是 K 的一个扩张。由于 A 是可分的，环 \(L \otimes_{K} A\) 是既约的；于是命题5表明环 \((L \otimes_{K} A) \otimes_{K} B\) （同构于 \(L \otimes_{K} (A \otimes_{K} B)\) ）是既约的，推论得证。

